\documentclass[11pt]{article}

\usepackage{amsmath,amssymb}
\usepackage[margin=1in]{geometry}
\usepackage{booktabs}
\usepackage{array}
\usepackage{hyperref}

\title{Short-Time Trace Asymptotic Feasibility of a Frozen Branch-Local Cochain-Laplacian Hierarchy}
\author{Tomislav Rupi\'c \\ Independent Researcher}
\date{March 2026}

\begin{document}

\maketitle

\begin{abstract}
Phase VIII of the frozen HAOS-IIP program tests one narrow question: whether the already frozen branch-local cochain-Laplacian hierarchy supports a reproducible short-time trace asymptotic regime across refinement. The phase does not modify Phase V, Phase VI, or Phase VII. It evaluates the exact trace proxy
\[
T_h(t) = \mathrm{Tr}\!\left(e^{-t \Delta_h}\right)
\]
on the deterministic hierarchy $h = 1 / n_{\mathrm{side}}$ with $n_{\mathrm{side}} \in \{12,24,36,48\}$. The trace is computed from the exact torus-mode spectrum induced by the frozen periodic DK2D block cochain-Laplacian structure, with direct dense validation on the coarsest level. A reproducible operational short-time window is identified on the fixed probe grid:
\[
t \in \{0.02, 0.03, 0.05, 0.075, 0.1\}.
\]
Across this window, the maximum cross-refinement log-slope dispersion is $0.005211$, the maximum normalized trace-spacing range is $0.005821$, the maximum kernel share is $0.010092$, and the maximum kernel-subtracted slope drift is $0.003044$. A size-factored quadratic surrogate
\[
T_h(t) \approx h^{-2}\bigl(b_0(h) + b_1(h)t + b_2(h)t^2\bigr)
\]
fits all refinement levels with $R^2 \ge 0.999996$, while the largest coefficient span across refinement is $0.027261$. Deterministic spectral cutoffs at $\lambda \le 7.5$ and $\lambda \le 7.8$ preserve the qualitative regime. The result is bounded: Phase VIII establishes short-time trace asymptotic feasibility for the frozen operator hierarchy. It does not establish continuum correspondence, geometric coefficients, universal asymptotic coefficients, or any physical law.
\end{abstract}

\section{Scope and Freeze Conditions}

Phase VIII is a feasibility instrument, not an interpretive layer. It consumes three already frozen authorities:
\begin{itemize}
\item the Phase V recovery-statistics authority manifest;
\item the Phase VI operator freeze manifest;
\item the Phase VII spectral-feasibility manifest.
\end{itemize}

The phase is not allowed to change the operator construction, the refinement hierarchy, the Phase V observable, or the Phase VII spectral baseline. Its only task is to determine whether a stable short-time trace regime can be identified and frozen on the current branch.

\section{Frozen Inputs}

The frozen operator is the branch-local periodic DK2D block cochain Laplacian
\[
\Delta_h,
\]
defined on the deterministic hierarchy
\[
h = \frac{1}{n_{\mathrm{side}}}, \qquad n_{\mathrm{side}} \in \{12,24,36,48\}.
\]

The authoritative references used by the phase are:
\begin{itemize}
\item \url{phase5-readout/phase5_authoritative_manifest.json};
\item \url{phase6-operator/phase6_operator_manifest.json};
\item \url{phase7-spectral/phase7_spectral_manifest.json};
\item \url{phase8-trace/phase8_trace_manifest.json}.
\end{itemize}

The master deterministic probe grid is fixed to
\[
\{0.01, 0.015, 0.02, 0.03, 0.05, 0.075, 0.1, 0.15, 0.2, 0.3, 0.5, 0.75, 1.0\}.
\]
This grid, together with the selected operational window, is now part of the frozen spectral-trace contract for later stages.

\section{Trace Proxy and Validation}

The trace proxy is defined by
\[
T_h(t) = \mathrm{Tr}\!\left(e^{-t \Delta_h}\right).
\]
No stochastic estimator is used. Instead, the trace is evaluated from the exact torus-mode spectrum induced by the block structure of the frozen periodic DK2D cochain-Laplacian. On the coarsest level, direct dense diagonalization and the exact mode-spectrum construction agree to machine precision, with maximum eigenvalue difference $0$ and maximum trace difference $0$ on the full probe grid.

Two deterministic cutoff probes are also recorded:
\begin{itemize}
\item $\lambda \le 7.5$;
\item $\lambda \le 7.8$.
\end{itemize}
These do not redefine the phase. They only test whether the qualitative short-time regime survives bounded spectral truncation.

\section{Operational Short-Time Window}

The phase searches the declared probe grid for the longest contiguous window satisfying simultaneous constraints on:
\begin{enumerate}
\item cross-refinement log-slope dispersion;
\item normalized trace-spacing stability;
\item kernel-share suppression;
\item truncation error bounds;
\item coefficient-fit stability.
\end{enumerate}

The selected operational short-time window is
\[
\{0.02, 0.03, 0.05, 0.075, 0.1\}.
\]

Its recorded worst-case diagnostics are:
\begin{center}
\begin{tabular}{@{}lr@{}}
\toprule
Metric & value \\
\midrule
max log-slope dispersion & 0.005211400534 \\
max normalized trace relative range & 0.005820971051 \\
max kernel share & 0.010092438116 \\
max relative trace error for $\lambda \le 7.5$ & 0.039036855531 \\
max relative trace error for $\lambda \le 7.8$ & 0.014977531175 \\
\bottomrule
\end{tabular}
\end{center}

The full trace remains monotone decreasing in $t$ at every refinement level.

\section{Trace and Slope Ledger}

\begin{table}[t]
\centering
\small
\begin{tabular}{@{}rrrrrr@{}}
\toprule
Level & $n_{\mathrm{side}}$ & $h$ & $T_h(0.02)$ & $T_h(0.075)$ & $T_h(0.1)$ \\
\midrule
R1 & 12 & 0.083333 & 532.859180 & 433.599978 & 396.336342 \\
R2 & 24 & 0.041667 & 2129.284951 & 1728.204937 & 1577.998557 \\
R3 & 36 & 0.027778 & 4789.988313 & 3885.867004 & 3547.422953 \\
R4 & 48 & 0.020833 & 8514.972194 & 6906.592208 & 6304.615453 \\
\bottomrule
\end{tabular}
\caption{Representative trace values on the frozen hierarchy inside the operational window.}
\label{tab:trace}
\end{table}

\begin{table}[t]
\centering
\small
\begin{tabular}{@{}rrrrrr@{}}
\toprule
Level & $S_h(0.02)$ & $S_h(0.03)$ & $S_h(0.05)$ & $S_h(0.075)$ & $S_h(0.1)$ \\
\midrule
R1 & -0.078531 & -0.118156 & -0.192277 & -0.277148 & -0.359390 \\
R2 & -0.079538 & -0.119652 & -0.194652 & -0.280474 & -0.363554 \\
R3 & -0.079726 & -0.119931 & -0.195095 & -0.281093 & -0.364329 \\
R4 & -0.079792 & -0.120029 & -0.195250 & -0.281310 & -0.364601 \\
\bottomrule
\end{tabular}
\caption{Log-trace slopes $S_h(t) = d(\log T_h)/d(\log t)$ across the selected window.}
\label{tab:slopes}
\end{table}

The slope table shows the main Phase VIII signal: the curves are not identical, but they remain tightly clustered and refinement ordered through the full selected window.

\section{Coefficient-Like Surrogate}

Phase VIII uses one descriptive surrogate only:
\[
T_h(t) \approx h^{-2}\bigl(b_0(h) + b_1(h)t + b_2(h)t^2\bigr).
\]
This is not interpreted as a geometric expansion. It is only a compact descriptive fit used to test drift and refinement smoothness.

\begin{table}[t]
\centering
\small
\begin{tabular}{@{}rrrrrr@{}}
\toprule
Level & $h$ & $b_0(h)$ & $b_1(h)$ & $b_2(h)$ & $R^2$ \\
\midrule
R1 & 0.083333 & 3.992541196696 & -15.196408636509 & 27.978374055231 & 0.999996635281 \\
R2 & 0.041667 & 3.992270826129 & -15.383158273830 & 28.598645942437 & 0.999996461490 \\
R3 & 0.027778 & 3.992219728715 & -15.417962259342 & 28.714851153143 & 0.999996428345 \\
R4 & 0.020833 & 3.992201767437 & -15.430159990250 & 28.755622542545 & 0.999996416672 \\
\bottomrule
\end{tabular}
\caption{Coefficient-like surrogate fits on the operational short-time window.}
\label{tab:coefficients}
\end{table}

The coefficient spans across refinement are:
\[
\mathrm{span}(b_0)=0.000085020802,\quad
\mathrm{span}(b_1)=0.015221236998,\quad
\mathrm{span}(b_2)=0.027260519706.
\]
The minimum fit quality is
\[
R^2_{\min} = 0.999996416672.
\]
The maximum probe-subgrid coefficient drift inside the selected window is $0.019376305038$.

\section{Truncation and Kernel Robustness}

The truncation robustness checks are:
\begin{center}
\begin{tabular}{@{}lrr@{}}
\toprule
Cutoff & max coefficient relative diff. & max log-slope diff. \\
\midrule
$\lambda \le 7.5$ & 0.110877321888 & 0.010940479377 \\
$\lambda \le 7.8$ & 0.044047599547 & 0.004232595036 \\
\bottomrule
\end{tabular}
\end{center}

The kernel contribution is also controlled inside the selected window. The maximum kernel share is $0.010092438116$, and the maximum log-slope difference between the full trace and the kernel-subtracted trace is $0.003044349031$.

These numbers justify a bounded feasibility statement. They do not justify a stronger asymptotic interpretation.

\section{Bounded Interpretation}

Phase VIII supports the following narrow authority statement:
\begin{quote}
The frozen branch-local cochain-Laplacian hierarchy supports a reproducible operational short-time trace regime on the declared probe grid, with stable slope behavior, smooth coefficient-like drift, and bounded truncation sensitivity across refinement.
\end{quote}

This statement is enough to freeze the spectral-trace baseline for later stages. It is not enough to identify geometric coefficients, universal asymptotic data, or a continuum object. Later phases may build on this baseline, but they may not silently redefine the operational window or the deterministic probe grid.

\section{Non-Claims}

This paper does not claim:
\begin{enumerate}
\item continuum correspondence;
\item geometric or universal asymptotic coefficients;
\item Seeley--DeWitt structure as an achieved result;
\item heat-kernel asymptotics beyond the present feasibility test;
\item physical law or correspondence.
\end{enumerate}

\section{Conclusion}

Phase VIII locks one bounded result on the frozen branch: a reproducible short-time trace regime exists on the declared probe grid, its slope behavior is stable across refinement, its coefficient-like quantities drift smoothly, and its qualitative structure survives deterministic truncation checks. This freezes the spectral-trace baseline for subsequent phases without granting any continuum or geometric interpretation.

Phase VIII establishes short-time trace asymptotic feasibility for the frozen operator hierarchy.

\end{document}
